\documentclass[11pt,a4paper]{article}
\usepackage[margin=1in]{geometry}
\usepackage{amsmath,amssymb,bm}
\usepackage{booktabs}
\usepackage{hyperref}

\newcommand{\E}{\mathbb E}
\newcommand{\erf}{\operatorname{erf}}

\title{The perpendicular term the mean-field theories drop: FCN2 against FCN3}
\author{}
\date{2026-11-10}

\begin{document}
\maketitle

The FCN2 theory with a per-neuron perpendicular norm $q$ and the $x_0$ Hermite tower to $\mathrm{He}_{11}$ matches the
exact posterior without any account of how a neuron's dependence on the perpendicular directions lines up with the
rest of the network. The FCN3 theory, built the same way, misses a group of hidden neurons at strong learnability whose
readout term comes half from exactly that. This note writes down the term both theories drop, shows where the two
networks differ in it, and says what is measured and what is not.

Notation: $x\sim\mathcal N(0,I_d)$, $x_0$ is the target direction and $x_\perp\in\mathbb R^{d-1}$ the other
coordinates; $\langle f|g\rangle=\E_x[f(x)g(x)]$. For a function $f(x)=f(x_0,x_\perp)$, $P_\parallel f$ is its average over
the perpendicular coordinates at fixed $x_0$,
\begin{equation*}
  (P_\parallel f)(x_0)=\int f(x_0,x_\perp)\,\mathcal N(x_\perp;0,I_{d-1})\,dx_\perp=\E\big[f\,\big|\,x_0\big]
  =\sum_{n\ge0}\langle\varphi_n|f\rangle\,\varphi_n(x_0),
\end{equation*}
with $\varphi_n=\mathrm{He}_n/\sqrt{n!}$: the orthogonal projection of $f$ onto the functions of $x_0$ alone, i.e.\ its
full $x_0$ Hermite expansion. For example $P_\parallel\erf(w\cdot x)=\erf\big(w_0x_0/\sqrt{1+2q}\big)$ with
$q=|w_\perp|^2$, since $w_\perp\cdot x_\perp$ is Gaussian with variance $q$. For a unit $k$ (a neuron of FCN2, a
second-layer hidden neuron of FCN3) with activation $g_k$, its perpendicular part is
$\delta g_k=g_k-P_\parallel g_k$, everything in $g_k$ that $P_\parallel$ removes; it is orthogonal to
every function of $x_0$.

\section*{The dropped term is the same in both networks}

Why the overlap with the residual matters. In the sampler the readout weights $a$ are integrated out exactly, and
removing unit $k$ from the network lowers the log-posterior by
$F_k=\frac{\sigma_A^2}{2N_2\chi}\,(g_k\cdot\alpha)^2/(1-\ell_k)$, plus a term $\frac12\log(1-\ell_k)$ that is
0.02--0.18 nats in size here and nearly the same for every unit, with $\alpha=(y-\bar f)/(\kappa/\chi)$ on the training
inputs, $\bar f=\sum_k\bar a_kg_k$, $\bar a$ the posterior mean of $a$, and $1-\ell_k$ the posterior variance of $a_k$ over
its prior variance. In the theory the same term is $\frac{\tilde g}{2}(\mathbf t\cdot\mathbf e)^2$ with
$\mathbf t=(y-\bar f)/(\kappa/P)$. In both, the unit enters through the overlap of its activation with the residual
$r=y-\bar f$: $g_k\cdot\alpha=\frac P\kappa\cdot\frac1P\sum_\mu g_k(x_\mu)r(x_\mu)$ in the sampler (taking $\chi=1$) and
$\mathbf t\cdot\mathbf e=\frac P\kappa\langle g_k|r\rangle$ in the theory. Below the overlap is taken over the population,
$\langle g_k|r\rangle$, as in the theory.

Splitting the overlap. The one fact needed is that a perpendicular part has zero overlap with any function of $x_0$
alone. Take any $f$ and any $u(x_0)$, and write the overlap as an integral over $x_0$ and then over $x_\perp$, with
$\mathcal N$ the standard normal densities. Neither $u(x_0)$ nor $(P_\parallel f)(x_0)$ depends on $x_\perp$, so both come
out of the inner integral, and $\int\mathcal N(x_\perp)\,dx_\perp=1$:
\begin{equation*}
  \langle f-P_\parallel f\,|\,u\rangle
  =\int\!\mathcal N(x_0)\,u(x_0)\left[\int\!\mathcal N(x_\perp)\,f(x_0,x_\perp)\,dx_\perp-(P_\parallel f)(x_0)\right]dx_0 .
\end{equation*}
The first term in the bracket is $(P_\parallel f)(x_0)$ by its definition, so the bracket is zero at every $x_0$ and the
overlap vanishes. Now write each activation as $g_k=P_\parallel g_k+\delta g_k$. $P_\parallel$ is an integral, so it is
linear, and the mean output splits the same way,
\begin{equation*}
  \bar f=P_\parallel\bar f+\bar f_\perp,\qquad
  P_\parallel\bar f=\sum_{k'}\bar a_{k'}P_\parallel g_{k'},\qquad
  \bar f_\perp=\sum_{k'}\bar a_{k'}\,\delta g_{k'},
\end{equation*}
and since the target depends on $x_0$ only ($P_\parallel y=y$), the residual is
$r=(y-P_\parallel\bar f)-\bar f_\perp$: a function of $x_0$ minus a perpendicular part. Expanding,
\begin{equation*}
  \langle g_k|r\rangle
  =\langle P_\parallel g_k\,|\,y-P_\parallel\bar f\rangle
  -\underbrace{\langle P_\parallel g_k\,|\,\bar f_\perp\rangle}_{=0}
  +\underbrace{\langle\delta g_k\,|\,y-P_\parallel\bar f\rangle}_{=0}
  -\langle\delta g_k\,|\,\bar f_\perp\rangle ,
\end{equation*}
where both cross terms vanish because one factor is a function of $x_0$ and the other a perpendicular part. This
leaves
\begin{equation}
  \langle g_k|r\rangle=\big\langle P_\parallel g_k\,\big|\,y-P_\parallel\bar f\big\rangle
  -\big\langle\delta g_k\,\big|\,\bar f_\perp\big\rangle,
  \qquad
  \bar f_\perp=\sum_{k'}\bar a_{k'}\,\delta g_{k'} .
  \label{eq:split}
\end{equation}
Both theories keep the first term, $\mathbf t\cdot\mathbf c(w_0,q)$ in FCN2 and $\mathbf t\cdot\mathbf e(\mathbf h,q_h)$ in
FCN3, and both set $\bar f_\perp=0$: their single-unit distribution does not prefer any perpendicular direction, so any
average of unit activations depends on $x_0$ only (apart from functions of $|x_\perp|$, which the FCN2 checks found
negligible). The dropped term is therefore
\begin{equation}
  \big\langle\delta g_k\big|\bar f_\perp\big\rangle=\sum_{k'}\bar a_{k'}\,Q_{kk'},
  \qquad Q_{kk'}=\langle\delta g_k|\delta g_{k'}\rangle ,
\end{equation}
and the two networks differ only in $Q_{kk'}$, the overlap of two units' perpendicular parts.

\section*{FCN2: units overlap only through $w_\perp\cdot w'_\perp$}

Here $g_k=\erf(w_k\cdot x)$ with $w_k=(w_{k0},w_{k\perp})$, $q_k=|w_{k\perp}|^2$ and $D_k=1+2|w_k|^2=1+2(w_{k0}^2+q_k)$.

First, $Q_{kk'}$ reduces to two overlaps of whole functions. Expanding the product,
\begin{equation*}
  Q_{kk'}=\langle g_k-P_\parallel g_k\,|\,g_{k'}-P_\parallel g_{k'}\rangle
  =\langle g_k|g_{k'}\rangle-\langle P_\parallel g_k|g_{k'}\rangle-\langle g_k|P_\parallel g_{k'}\rangle
  +\langle P_\parallel g_k|P_\parallel g_{k'}\rangle .
\end{equation*}
In the second term write $g_{k'}=P_\parallel g_{k'}+\delta g_{k'}$:
$\langle P_\parallel g_k|g_{k'}\rangle=\langle P_\parallel g_k|P_\parallel g_{k'}\rangle+\langle P_\parallel g_k|\delta g_{k'}\rangle$,
and the last overlap is zero (a function of $x_0$ against a perpendicular part). The third term is the same with $k$
and $k'$ exchanged. So both middle terms equal $\langle P_\parallel g_k|P_\parallel g_{k'}\rangle$, and
\begin{equation*}
  Q_{kk'}=\langle g_k|g_{k'}\rangle-\langle P_\parallel g_k|P_\parallel g_{k'}\rangle .
\end{equation*}

Both overlaps follow from one identity: for zero-mean jointly Gaussian $z,z'$,
\begin{equation*}
  \E\big[\erf(z)\erf(z')\big]=\frac2\pi\arcsin\frac{2\,\mathrm{cov}(z,z')}{\sqrt{(1+2\,\mathrm{var}\,z)(1+2\,\mathrm{var}\,z')}} .
\end{equation*}
For the whole activations take $z=w_k\cdot x$, $z'=w_{k'}\cdot x$ with $x\sim\mathcal N(0,I_d)$:
$\mathrm{var}\,z=|w_k|^2$, so $1+2\,\mathrm{var}\,z=D_k$, and $\mathrm{cov}(z,z')=w_k\cdot w_{k'}=w_{k0}w_{k'0}+u_{kk'}$ with
$u_{kk'}=w_{k\perp}\cdot w_{k'\perp}$. Hence
\begin{equation*}
  \langle g_k|g_{k'}\rangle=\frac2\pi\arcsin\frac{2(w_{k0}w_{k'0}+u_{kk'})}{\sqrt{D_kD_{k'}}} .
\end{equation*}
For the $x_0$ parts, $P_\parallel g_k=\erf(s_kx_0)$ with $s_k=w_{k0}/\sqrt{1+2q_k}$ (the example in the notation
paragraph). Take $z=s_kx_0$, $z'=s_{k'}x_0$: $\mathrm{cov}=s_ks_{k'}$ and
$1+2s_k^2=(1+2q_k+2w_{k0}^2)/(1+2q_k)=D_k/(1+2q_k)$, so
\begin{gather*}
  \frac{2s_ks_{k'}}{\sqrt{(1+2s_k^2)(1+2s_{k'}^2)}}
  =\frac{2w_{k0}w_{k'0}}{\sqrt{(1+2q_k)(1+2q_{k'})}}\sqrt{\frac{(1+2q_k)(1+2q_{k'})}{D_kD_{k'}}}
  =\frac{2w_{k0}w_{k'0}}{\sqrt{D_kD_{k'}}},\\
  \langle P_\parallel g_k|P_\parallel g_{k'}\rangle=\frac2\pi\arcsin\frac{2w_{k0}w_{k'0}}{\sqrt{D_kD_{k'}}} .
\end{gather*}
Subtracting the two,
\begin{equation}
  Q_{kk'}=\frac{2}{\pi}\left[\arcsin\frac{2(w_{k0}w_{k'0}+u_{kk'})}{\sqrt{D_kD_{k'}}}
  -\arcsin\frac{2w_{k0}w_{k'0}}{\sqrt{D_kD_{k'}}}\right].
  \label{eq:Q2}
\end{equation}
The perpendicular weights of the two units enter only through $u_{kk'}$, and $Q_{kk'}=0$ when $u_{kk'}=0$.

Which part of the activation each power of $u$ belongs to. At fixed $x_0$, $\xi=w_{k\perp}\cdot x_\perp$ and
$\xi'=w_{k'\perp}\cdot x_\perp$ are jointly Gaussian with covariance $u_{kk'}$. Mehler's formula expands the average of a
product of functions of $\xi$ and $\xi'$ in powers of that covariance, and the $u^m$ term is the overlap of the two
functions' degree-$m$ Hermite parts in $\xi$ and $\xi'$; so Taylor-expanding \eqref{eq:Q2} in $u$,
$Q_{kk'}=\sum_{m\ge1}b_m\,u_{kk'}^m$, the $u^m$ term is the overlap of the two units' parts of degree $m$ in $x_\perp$. A
unit's degree-$m$ perpendicular part points along $\hat w_{k\perp}^{\otimes m}$, a direction that belongs to that unit
alone. For independent directions $u\sim\sqrt{q_kq_{k'}/(d-1)}$, so the degree-$m$ overlap between different units falls
as $d^{-m/2}$, whatever the units do. Only the linear part ($m=1$, the $d-1$ directions of $x_\perp$) is shared between
units.

\section*{FCN3: hidden units are built from one shared set of functions}

The hidden pre-activation splits exactly into
\begin{equation}
  h_k(x)=a_k(x_0)+h_{k\perp}(x),\qquad h_{k\perp}=\sum_jB_{kj}\,\delta\psi_j,\qquad
  \delta\psi_j=\erf(w_j\cdot x)-P_\parallel\erf(w_j\cdot x),
\end{equation}
so every hidden unit's perpendicular field is a combination, with its own Gaussian weights $B_k$, of the same $N_1$
read-in functions $\delta\psi_j$. Treating $h_{k\perp}$ as Gaussian over $x_\perp$ at fixed $x_0$, as (E2) does, the part
of $\delta g_k$ linear in $h_{k\perp}$ is (Stein's lemma)
\begin{equation}
  \delta g_k=\gamma_k(x_0)\,h_{k\perp}+\dots,\qquad
  \gamma_k(x_0)=\E_\perp\big[\erf'(h_k(x))\,\big|\,x_0\big]
  =\frac{2}{\sqrt\pi}\,\frac{e^{-a_k(x_0)^2/(1+2q_{h,k})}}{\sqrt{1+2q_{h,k}}},
\end{equation}
where the dots are of second and higher order in $h_{k\perp}$ and orthogonal to the first. At this order
\begin{gather}
  Q_{kk'}\simeq\sum_{i,j}B_{ki}\,\big\langle\gamma_k\delta\psi_i\big|\gamma_{k'}\delta\psi_j\big\rangle\,B_{k'j},\\
  \big\langle\delta g_k\big|\bar f_\perp\big\rangle\simeq B_k\cdot\mathbf U,\qquad
  U_i=\sum_{k'}\bar a_{k'}\sum_j\big\langle\gamma_k\delta\psi_i\big|\gamma_{k'}\delta\psi_j\big\rangle B_{k'j}.
\end{gather}
The dropped term is a linear function of the unit's own Gaussian weights. (E2) keeps two functions of $B_k$, its
$x_0$ projections $\mathbf h_k=CB_k$ and the size of its perpendicular field $q_h=B_k^\top AB_k$
($A_{ij}=\langle\delta\psi_i|\delta\psi_j\rangle$), and averages $B_k\cdot\mathbf U$ away.

\section*{Where the networks differ}

In FCN3 the diagonal entries $\langle\delta\psi_j|\delta\psi_j\rangle=P_\perp(w_j)$ are of order one, and every hidden
unit uses every read-in neuron's whole perpendicular part, its higher-degree part included. Two hidden units' parts of
degree $m\ge2$ in $x_\perp$ therefore overlap through $\sum_jB_{kj}B_{k'j}\,\|\delta\psi_j^{(m)}\|^2$: of order
$N_1^{-1/2}$ for independent weight rows, of order one for a group of rows that point the same way. In FCN2 the same
overlaps fall as $d^{-m/2}$ and no choice of parameters changes that. The output's perpendicular part
$\bar f_\perp\simeq\gamma\sum_j\beta_j\delta\psi_j$, $\beta=\sum_{k'}\bar a_{k'}B_{k'}$, lies in the span that every hidden
unit draws on.

At linear order, however, the two networks look alike. An FCN2 unit has $d-1=149$ perpendicular directions with
variance $q/(d-1)\approx0.0034$ each; an FCN3 hidden unit's perpendicular field has an effective number of directions
$M=(\operatorname{tr}G_\perp)^2/\operatorname{tr}G_\perp^2\approx141$--164 (theory 164; 141--148 from the exact prior of the
saved networks) and variance $\bar q_h/M\approx0.003$ each. In both a unit can turn its perpendicular part towards
$\bar f_\perp$ at a Gaussian prior cost. So the equations alone do not say that FCN3 needs the dropped term and FCN2
does not; that depends on how large $\bar f_\perp$ is in each network and whether the shared higher-degree channel
carries it.

\section*{What is measured}

FCN3, onset scan at $d=150$, $N_1=N_2=1500$, $\kappa=0.015$, $P=48000$, last saved network of each of the three
datasets. In every network the same 17 of the 1500 hidden units have $|h_3|>0.3$ for the whole run (the large-$h_3$
group the theory does not have). On 48000 fresh inputs, with $\delta g_k$ computed from $g_k$ minus its
$x_0$ Hermite part up to $\mathrm{He}_{21}$ (using (E2)'s $\mathbf e$ at the unit's own $\mathbf h$, $q_h$):
\begin{center}
\begin{tabular}{@{}lc@{}}
\toprule
$|\cos|$ between two units' $\delta g$: ordinary pairs & 0.065\\
\phantom{$|\cos|$ between two units' $\delta g$:} pairs within the 17 & 0.08--0.10\\
rms of $\bar f_\perp$ from all units except the 17 & 0.16--0.17\\
rms of $\bar f_\perp$ from the 17 alone & 0.10--0.11\\
cosine between those two parts & $-0.81$ to $-0.83$\\
rms of the total $\bar f_\perp$ & 0.10\\
\bottomrule
\end{tabular}
\end{center}
The output's perpendicular part is small next to the output and not small next to its error. Fitting the mean output on
200000 fresh inputs by functions of $x_0$ alone (Hermite degrees 0 to 31) and taking the remainder gives, in the three
networks, a mean square of $\bar f_\perp$ of 0.0088--0.0097 against 1.47 for the whole output (0.6\%), while the test
error $\|y-\bar f\|^2=0.032$--0.035 splits into 0.022--0.026 from the $x_0$ part and those 0.009 from $\bar f_\perp$
(about 28\%). The theory sets that last part to zero.

Pairwise the 17 are nearly independent; together they cancel a large part of what the other units put into
$\bar f_\perp$ (mean square 0.026--0.029 from the other units, 0.010 after the 17 are added). Their readout
term in the sampler is 470--535 nats each against 225--250 for the theory's term at the same $\mathbf h$, $q_h$, and
about half of the sampler's comes through $\langle\delta g_k|\bar f_\perp\rangle$, which is as large on fresh data
as on the training set.

FCN2: not measured. The runs on the same line ($d=150$, $N=1500$, $\kappa_0=0.015$, $P=3000$--24000) are saved in
\texttt{milestones/fcn2\_collapsed\_hmc/runs/*pscan\_d150*}. Measuring $\bar f_\perp$ and each unit's
$\langle\delta g_k|\bar f_\perp\rangle$ there decides between the two readings: a comparable $\bar f_\perp$ that no
group cancels points to the shared higher-degree channel of FCN3; a small $\bar f_\perp$ points to the strength of
learning (FCN2 reached $L_3\approx0.38$ on this line, FCN3 $0.90$ at $P=48000$) or to something else.

\end{document}
